\paragraph{Finite splitting check.}
The checker counts the roots of all six quintics over each good prime field through 5000, omitting $2,5$ and the prime divisors of the quadratic and polynomial discriminants. For $D=-12451$ it checks 663 primes. The first with all six splitting is 269; at 11 only the quintic with $p=7$ in the table splits completely. For $D=-37363$ it checks 661 primes and the first common splitting prime is 769. The first representative of every splitting case is also factored independently. The algebraic argument gives the densities, and the finite data illustrate the distinctions on which the analytic test depends.

\section{Further research questions}
\begin{enumerate}
\item Can the exponent-seven search be closed below $4\cdot4600^7$? For seventh powers, $q^7-p^7$ has a different factorization and a larger bound for the repeated coefficient. A usable extension would need to prove that two small split primes yield compatible norm equations or find a replacement elimination. This is the one exponent among $1,\ldots,8$ whose global census remains unresolved here.
\item For prime-gap graphs, what is the asymptotic minimum of the largest matching among all realizations on $n$ vertices? The block theorem gives a lower bound of order $p_n/(\log p_n)^2$; a better bound lengthens the block of insertions that use disjoint old endpoints.
\item For $D=-12451$ the unique tame line gives an everywhere-finite-unramified cyclic quintic extension of the quartic mirror field. Can a theta-lift or height construction estimate its splitting primes with the sign and square-root accuracy of Section~\ref{sec:joint-test}, first for this line alone? The other five lines have a known conductor at five, so this family distinguishes that first case from its ramified companions.
\end{enumerate}
